\begin{EnglishAbstract}
	
	``
	Duchenne muscular dystrophy (DMD) is a progressive neuromuscular disorder characterized by gradual skeletal-muscle weakness and declining motor function. The North Star Ambulatory Assessment (NSAA) is commonly used to evaluate ambulatory motor function in individuals with DMD and assigns a score of 0, 1, or 2 to each task based on clinical observation. Because this scoring process partly depends on the evaluator’s judgment, quantitative and repeatable methods for analyzing task performance may provide useful complementary information. This study investigated the technical feasibility of classifying simulated functional levels corresponding to selected NSAA item scores using wearable inertial sensors and machine-learning algorithms.
	No motion data were collected directly from patients with DMD. The development dataset was generated by one healthy participant familiar with DMD-related motor limitations and the NSAA scoring criteria, who simulated three functional performance levels under controlled conditions. Five activities were investigated: standing, rising from a chair, two-footed jumping, lifting the head from a supine position, and stepping onto a box with the right foot. Five inertial measurement units were attached to the forehead, both wrists, and both ankles to record triaxial linear acceleration and angular velocity at approximately 90 Hz. Signal processing included fourth-order zero-phase Butterworth low-pass filtering with a 4-Hz cutoff frequency, and signal preparation, with activity segmentation additionally evaluated where applicable. Twenty-three time-domain features were extracted from each of the 30 signal channels, resulting in 690 features per trial. Standardization and principal component analysis (PCA) were implemented within the model-training pipeline to prevent information leakage. Support vector machine (SVM), k-nearest neighbors, and random forest classifiers were compared using repeated stratified five-fold cross-validation.
	The RBF-kernel SVM combined with PCA provided the most suitable performance for all five activities. Depending on the activity, 10–15 principal components were selected to obtain stable classification performance. The selected models met the predefined internal accuracy threshold of 90\% for standing, jumping, and head lifting, and 80\% for chair rise and right-foot step-up. These thresholds were engineering criteria used for model selection and should not be interpreted as measures of clinical validity.
	The findings demonstrate the technical feasibility of distinguishing simulated functional patterns in data collected from a single healthy participant. They do not establish the system’s performance in patients with DMD or its generalizability to new individuals. Clinical validation will require data collection from patients, simultaneous scoring by qualified assessors, and subject-independent evaluation.
	``
	
	\enkeywords{
		Duchenne muscular dystrophy, North Star Ambulatory Assessment, inertial measurement unit, motion analysis, machine learning, principal component analysis, motor function assessment
	}
	
\end{EnglishAbstract}